\chapter{Conclusions and Future Work}
\label{cap:conclusions}

\section{Conclusions}
This work has been carried out within the framework of the European and Ibero-American project \textit{EA-DIGIFOLK}, with the purpose of designing, implementing, and evaluating an automated, robust, and scalable architecture capable of overcoming traditional methodological limitations in the analysis of large folk musical corpora. To address the problems of format heterogeneity and the context limitations of general-purpose language models (LLMs), an intermediate infrastructure was proposed based on the extraction and normalization of metadata using specialized tools like \textit{music21} and storing the information in a structured relational database (\textit{SQLite}).

The evaluation and comparative analysis of this system against alternative approaches have yielded decisive results regarding the behavior of artificial intelligence applied to computational musicology. The study demonstrates that direct access to a structured database through local SQL inference (utilizing \textit{Llama3.1:8B}) offers a substantially more deterministic, precise, and consistent data extraction than systems based purely on file retrieval and search (\textit{Retrieval-Augmented Generation} or RAG with large and general models like \textit{Gemini 2.5 Flash}). While the RAG approach stands out for its linguistic flexibility and synthesis capacity, its propensity for hallucinations and generating content not backed by the dataset critically diminishes its reliability in specialized research settings, such as technical musicology. Nonetheless, the results also evidence that both approaches suffer when faced with queries that require complex, multi-step musical reasoning.

Based on the findings and the detailed discussion, the following main conclusions are drawn:

\begin{itemize}
    \item Structured storage in SQLite combined with a local model allows for highly reliable and hallucination-free information retrieval when the required data is explicitly represented in the schema. Conversely, RAG systems based on File Search provide greater contextual and descriptive richness in natural language, but introduce inaccuracies and false positives (such as inventing pieces that do not exist in the corpus).
    \item The primary limitation of the structured local approach does not reside in the availability of analytical data, but rather in the model's capacity to correctly translate complex natural language questions into executable SQL queries. Errors in this intermediate translation constitute a critical failure source that decreases the system's robustness against non-trivial or compound queries.
    \item Neither of the two evaluated approaches is fully autonomous or robust when directly resolving queries that involve complex abstractions or multi-step mathematical and musical calculations (such as analyzing boundary intervals, syncopations, or advanced melodic structures). Both systems tend to simplify responses or assume unjustified correlations in the absence of a prior analytical processing layer.
    \item For LLM-based tools to be truly useful in the musicological research of large collections, data representation must not only be structured but explicitly optimized for analytical queries. It is concluded that to improve these tools and their results, designing intermediate architectures that hybridize guided query generation, strict result validation, and specialized structured reasoning modules is necessary.
\end{itemize}

\section{Future Work}
The developed tool constitutes a first step toward applying artificial intelligence and computational musical analysis techniques to digital score corpora. Nonetheless, several lines of research could be explored in future work to expand the analytical capabilities of the tool, improve response quality, and enrich user experience.

A first line of development consists of implementing multi-agent architectures based on the \textit{Agent-to-Agent} (A2A) protocol. Instead of delegating the entire reasoning process to a single conversational agent, tasks could be distributed among agents specialized in different areas, such as musical analysis, information retrieval, musicological interpretation, or response generation. This approach would allow complex queries to be addressed through collaborative processes among agents, facilitating system specialization and functional scalability.

Another possible evolution involves incorporating \textit{Graph Retrieval-Augmented Generation} (GraphRAG) techniques. While the current implementation stores metadata in a relational database, a graph-based model would allow for the explicit representation of relationships existing among songs, authors, genres, geographical regions, musical structures, and stylistic features. This representation would facilitate more complex queries based on semantic and musicological connections, as well as the automatic exploration of non-evident relationships within the corpus. Likewise, the integration of hybrid models combining the current database with graph-based knowledge structures could be investigated. In this way, it would be possible to maintain the efficiency of structured queries while simultaneously taking advantage of relation-based contextual retrieval systems.

Another prospective line of research involves implementing musical similarity search mechanisms based on vector representations (\textit{embeddings}). From the metadata extracted during preprocessing, each work could be represented by a numerical vector summarizing its main musical and textual characteristics, such as key, mode, time signature, rhythmic patterns, metric phenomena, or lyric themes. This representation would allow for the automatic calculation of the degree of similarity between songs and enable new features, such as searching for related works, identifying repertoires with common traits, or automatically recommending songs.

On the other hand, a priority technical development line to improve the local model's precision lies in optimizing relational queries over encapsulated semi-structured data. Currently, the \textit{SQLite} database schema stores musical variables of high relational density and variability in text columns under structured JSON format; this is the case for critical fields such as \texttt{mapeo\_silabas\_duraciones} (which maps the exact correspondence between lyric syllables and musical durations) or \texttt{accidentales\_compases\_json} (which records the measure location of accidental alterations like \textit{F\#5} or \textit{C\#5}). With the objective of eliminating the multi-step reasoning limitations identified in LLMs, an immediate future task will consist of empowering the agent to execute native extraction functions like \texttt{json\_extract()}. Integrating and training the model on the syntax of these functions will allow it to isolate, query, and operate directly on the internal keys and values of JSON objects without needing to alter the rigid relational schema or overload the model's context with prior natural language destructuring. This would translate into an increase in system robustness when handling complex analytical questions that cross, for example, specific syncopations with syllable durations or accidental alterations in designated measures.

Lastly, a user-oriented enhancement would involve developing a specific interface for song comparison. This functionality would allow selecting two or more works and jointly visualizing their musical, textual, and structural characteristics. Among other aspects, keys, modes, time signatures, rhythmic patterns, event density, syncopation phenomena, polyrhythms, hemiolas, or textual themes could be compared. A tool of this type would facilitate comparative studies and could be especially useful in contexts of musicological research and traditional repertoire analysis.

In addition to these concrete lines, the modular nature of the proposed architecture allows for the incorporation of new metadata extraction algorithms as new analytical needs are identified. This opens the door to future investigations focused on the automatic detection of more complex musical phenomena and the progressive expansion of the knowledge base used by the system.